%! AP Statistics — Unit 2, Lesson 2.3 — Exit Ticket (Answer Key)
\documentclass[10pt]{article}
\usepackage{apstats-article}
\usepackage{apstats-key}

%=============================================================================
\begin{document}

\pageheader{Unit 2, Lesson 2.3}{Exit Ticket --- Answer Key}
\begin{tabularx}{\linewidth}{X X X}
Name:\;\ans{Key} & Date:\; & Period:\;
\end{tabularx}

\vspace{0.22in}

\begin{tcolorbox}[colback=sky, colframe=navy, boxrule=0.9pt, arc=2mm,
    left=4mm, right=4mm, top=3mm, bottom=3mm]
\small\textit{A health survey classified 330 adults by smoking status and whether they
developed a respiratory condition. Of 150 smokers, 45 developed the condition.
Of 180 non-smokers, 18 developed the condition.}
\end{tcolorbox}

\vspace{0.2in}

\begin{enumerate}[leftmargin=1.6em, itemsep=0.2in]

    \item $\hat{p}_{\text{smoker}} = \dfrac{\ans{45}}{\ans{150}} = \ans{0.30 = 30\%}$
          \hfill
          $\hat{p}_{\text{non-smoker}} = \dfrac{\ans{18}}{\ans{180}} = \ans{0.10 = 10\%}$

    \item $\hat{p}_{\text{smoker}} - \hat{p}_{\text{non-smoker}} = \ans{0.30} - \ans{0.10} = \ans{0.20}$\\[8pt]
          \ansline{The proportion of smokers who developed the condition (30\%) was 20 percentage
          points higher than the proportion of non-smokers who developed the condition (10\%).}

    \item $\text{RR} = \dfrac{\ans{0.30}}{\ans{0.10}} = \ans{3.0}$\\[8pt]
          \ansline{Smokers were 3.0 times as likely to develop the respiratory condition as
          non-smokers.}

    \item \ansline{Both statistics reveal a strong association. The relative risk (3.0) is
          particularly striking --- ``3 times as likely'' is easy for a general audience to
          understand. The absolute difference (20 pp) is also large enough to be clinically
          meaningful, so reporting both gives the most complete picture.
          (Accept any well-reasoned justification.)}

\end{enumerate}

\vspace{0.18in}

\begin{tcolorbox}[colback=redbg, colframe=redacc, boxrule=0.8pt, arc=2mm,
    left=4mm, right=4mm, top=3mm, bottom=3mm]
\small\textbf{AP-Style Practice}\\[6pt]
RR $= 3.0$ and difference $= 0.02$.

\begin{enumerate}[label=\textbf{(\Alph*)}, leftmargin=2em, itemsep=6pt]
  \item $\hat{p}_1 = 0.06$, $\hat{p}_2 = 0.04$: RR $= 1.5$, diff $= 0.02$. RR wrong.
  \item $\hat{p}_1 = 0.60$, $\hat{p}_2 = 0.20$: RR $= 3.0$, diff $= 0.40$. Diff wrong.
  \item $\hat{p}_1 = 0.03$, $\hat{p}_2 = 0.01$: RR $= 3.0$, diff $= 0.02$.
        \textcolor{keyred}{\textbf{$\leftarrow$ correct}}
  \item $\hat{p}_1 = 0.30$, $\hat{p}_2 = 0.10$: RR $= 3.0$, diff $= 0.20$. Diff wrong.
\end{enumerate}

\vspace{6pt}
\ans{(C). Only $\hat{p}_1 = 0.03$ and $\hat{p}_2 = 0.01$ satisfy \emph{both} conditions:
RR $= 0.03/0.01 = 3.0$ and difference $= 0.03 - 0.01 = 0.02$. This illustrates how a rare
outcome can have a large relative risk but a tiny absolute difference --- a core concept of today's lesson.}
\end{tcolorbox}

\begin{teachernote}
\textbf{Item 4 (open-ended):} Award full credit for any answer that correctly references
both statistics and states which one the student finds more compelling with a valid reason.
The AP exam often asks ``which better communicates'' --- both interpretations can earn credit
if reasoning is sound.

\textbf{AP item (C):} Distractors (B) and (D) have the correct relative risk but the wrong
difference. Distractor (A) has the correct difference but wrong RR. Students who confuse the
two statistics will likely choose (A) or (D). Use to reinforce that RR and the absolute
difference are independent properties of the data.
\end{teachernote}

\end{document}
